% =====================================================================
\section{No-gos: what each class cannot do}
\label{sec:nogos}
% =====================================================================

A no-go is a theorem that a \emph{class} of approaches cannot succeed, not a statement that no
approach can. The campaign proved several, and they are as valuable as the theorems because they
close branches. This chapter states each with the precise reason it fails. Where a no-go is a
theorem about a class, it is labelled \LN; where the underlying statement is a genuine theorem
about the theory (for example, that zero-temperature pure Yang--Mills is a single phase), the
underlying statement is labelled \LT{} and the no-go is its corollary.

\subsection{The Schur-convex master obstruction}
\label{sec:schur-master}

\textbf{Statement.} Any $g$-continuous, protected functional of the measure that vanishes at
$g=0$ (the free field) cannot produce a positive atom at $g>0$ by positivity or analyticity
alone. The free field satisfies all OS axioms and has $f_1^2=0$; any positivity-only argument
forcing $f_1^2>0$ would force it there too.

\textbf{Precise form.} Let $\mathcal F$ be the class of functionals $\Phi(\nu)$ that are (i)
continuous in the weak-$*$ topology restricted to finitely many moments, (ii) non-negative on
every OS measure, and (iii) vanish at the free-field measure $\nu_0$ (supported on $[0,\infty)$
with $f_1^2=0$). Then for every $\nu$ in the feasibility set and every $\varepsilon>0$ there is
$\nu'$ in the feasibility set with $f_1^2(\nu')=0$ and $|\Phi(\nu)-\Phi(\nu')|<\varepsilon$.
Hence no $\Phi\in\mathcal F$ can imply $f_1^2>0$. The proof is the smoothed-replacement
construction of Chapter~\ref{sec:schur}: the atomless bump matches the finite data.

\textbf{Consequence.} Positivity, reflection positivity, spectral positivity, and analyticity
give \textbf{upper bounds and the vacuum}, never the edge. This is the single most important
no-go; it is why the campaign always says ``positivity is blind to the edge.'' It is also the
one no-go that is not specific to Yang--Mills: it holds for every OS theory, which is exactly
why it cannot carry YM-specific information.

\subsection{Goldstone uniqueness}

The only positivity-plus-analyticity theorem that forces an atom is Goldstone's theorem, and it
lives at \textbf{zero mass} (broken continuous symmetry). The atom $f_1^2$ is the ``positive-mass
cousin'' of a Goldstone mode, and the positive-mass cousin does not exist at the level of a
general theorem. The one protected positive-mass atom is BPS, but it is in the \textbf{wrong
channel} (central charge/flux; $Z=0$ on $\operatorname{tr}F^2$). So no symmetry theorem supplies
a positive-mass atom in the scalar channel.

\subsection{The periodic no-go (no convexification)}
\label{sec:periodic}

\textbf{Statement.} The log-density of $\mu$ is periodic on the compact group; a periodic
concave function is constant, so there is no log-concave substitute equal to $\mu$. The Savvidy
negativity is the cost of periodic winding at $|x|>\pi/2$. Thermal smoothing gives only
semi-concavity, and the Pr\'ekopa lift is blocked by the indefinite $G''/H_{ff}$.

\textbf{Consequence.} Convexification of the effective action is impossible; the wall is real,
not a gauge artifact. The measured $7.1\%$ negative directions after the gauge quotient (W4)
are the quantitative form of the statement. Note carefully what this no-go does \emph{not} say:
it does not say the true vacuum is unstable (the Savvidy negative mode is a barrier-top artifact,
Chapter~\ref{sec:savvidy}); it says that no globally log-concave proxy exists on the compact
group. The two statements are compatible, and distinguishing them was one of the Savvidy wave's
contributions.

\subsection{The nested-cone no-go}

\textbf{Statement.} Every sign-changing / Karlin cone $K_0\subset K_1\subset\cdots$ contains the
positive cone, so on each layer it swallows the strictly positive free field and is vacuous. A
cone that passes must cut \emph{below} positivity (where the symbol is indefinite), and there it
presupposes the very atom it is meant to prove (circular).

\textbf{Consequence.} The ``second cone'' route to the edge is closed. Plaquette products do not
tensorize to a sign-regular cone. This is the structural reason the gauge-invariant
Perron--Frobenius theorem reaches the vacuum but not the glueball
(Chapter~\ref{sec:theorems}): the quotient has no order, so a Karlin flag cannot be built.

\subsection{The rigidity no-go (no index)}
\label{sec:rigidity}

\textbf{Statement.} Any index that is a discrete integer and locally constant under
symmetry-preserving deformations is constant along the connected $\beta$ family (which reaches
the free Gaussian at $\beta=\infty$), hence equals its free value, hence is $0$ for YM. The gap
must be a \textbf{$\beta$-non-analytic jump, not any discrete invariant.}

\begin{claimbox}{cN}
\claimhead{Proposition}{(rigidity)} \emph{Let $I(\beta)\in\mathbb Z$ be locally constant on
$(0,\infty)$ and defined by a homotopy invariant of the gauge-field configuration space. Then
$I(\beta)=I(\infty)=0$ for every $\beta$ along the Wilson family.}
\begin{proofx}
The Wilson family $\beta\mapsto\mu_\beta$ is continuous in $\beta$ and connects every finite
$\beta$ to the Gaussian limit $\beta\to\infty$ (in the sense of the transfer operator's
spectrum). A homotopy invariant is unchanged by continuous deformation; in the Gaussian limit the
only index available is the trivial one. Hence $I$ is constant and equals its free value. The
only acid-passing index in the corpus --- the Gribov--Faddeev--Popov count --- counts gauge
copies, not physical states, and is not an invariant of $\mu_\beta$.
\end{proofx}
\end{claimbox}

\textbf{Consequence.} Noncommutative geometry, localization, BPS, and Witten-spinor indices all
return $0$ at the free point; none can carry the gap. The statement is robust: it applies to any
discrete invariant, not to a particular construction. This is why the campaign says the gap is
``not topological''; the reason is not a failure of a construction but the constancy of indices.

\subsection{The Feshbach / Schur angle no-go}

\textbf{Statement.} Block-decompose $H=\begin{psmallmatrix}A&B\\B^{\mathsf T}&C\end{psmallmatrix}$
with $C\prec0$. Then $\Sigma=A-BC^{-1}B^{\mathsf T}=A+B|C|^{-1}B^{\mathsf T}\succeq A$:
eliminating a negative block \textbf{only adds} to the survivor (Haynsworth inertia; numerically
verified $\lambda_{\min}(\Sigma-A)=-6\times10^{-16}$). So ``a small angle preserves $\Sigma>0$''
is a sign error: $\Sigma>0\iff A=P^{\mathsf T}HP>0$, and $A$ \emph{is} the gap. The only way an
angle enters is through the overlap $\varepsilon=\sin^2\theta$, requiring
$\varepsilon<\tfrac12$ ($\theta<45^\circ$); on a frozen constant background $\theta=0$ but the
statement is vacuous ($F\equiv0$, $g$-dependent), and on the physical $k=0$ mode
$\varepsilon=O(1)$ (measured $P^{\mathsf T}HP$ minimum $-0.3$) with no momentum separation. The
whole route is $\iff\varepsilon<\tfrac12\iff$ H1.

\subsection{The Krein / Pontryagin no-go}

\textbf{Statement.} Krein counting/sign invariants live at the UV scale or count signs; the gap
is the IR \textbf{min}. The two are mutually exclusive, as detailed in Chapter~\ref{ch:tools}.
The residual is H1. The no-go is a statement about the \emph{type} of invariant: a count is
rigid, the gap is a threshold; a UV scale cannot see an IR scale.

\subsection{The order-parameter no-go (single-phase theorem)}
\label{sec:b2}

This is the deepest no-go of the roof wave, and its underlying statement is a genuine theorem.

\begin{claimbox}{cT}
\claimhead{Theorem}{(B2: single phase)} \emph{Zero-temperature four-dimensional pure
Yang--Mills is a single phase: the $\beta$ dependence is a smooth crossover, not a phase
transition. Consequently there is no order parameter and no defect-condensation criterion in
any channel that could certify the $0^{++}$ gap.}
\end{claimbox}

The proof is the corpus's (\code{PROOF-NO-MELT-PHASE.md}, \code{ROOF-TRANSFER.md}): the
$T=0$ system has no local order parameter; the finite-temperature deconfinement transition at
$\beta_c(N_t)$ is a property of the thermal circle, not of the zero-temperature vacuum; and the
crossover observed on the lattice is smooth in $\beta$. The corollary is the no-go: an
order-parameter criterion needs a transition to condense, and a single phase has none. The 3D
$Z_2$ specimen models \emph{confinement} (a transition), not the single-phase $0^{++}$ gap. This
eliminates the entire order-parameter / defect-condensation paradigm for the roof, independently
of the channel-mismatch arguments.

\subsection{The topological and Markov eliminations}

\textbf{Topology.} Dead by rigidity (\S\ref{sec:rigidity}): any index is constant on the
connected $\beta$ family and counts the atom side, certifying $\text{gap}=0$ rather than
$\text{gap}>0$. The statement is stronger than a failure of a particular construction: it is a
theorem about all discrete invariants.

\textbf{Markov-semigroup / gradient-flow functionals.} Dead by Theorem BE
(\code{BASE-FIELDMEASURE.md}): a functional constructed from the semigroup of the transfer
operator is homogeneous in the Hessian (so it is $\kappa$ relabelled), zero at equilibrium, or
sign-indefinite. In all three cases it cannot produce a lower bound on $G''$ that is positive at
the free point and negative in the gapless direction. The reason is structural: the semigroup is
the same object whose gap is being measured, so a functional of it either presupposes the gap or
cancels at equilibrium.

\subsection{The analytic-transfer no-go (essential singularity)}
\label{sec:esssing}

The gap is
\begin{equation}
m\ \sim\ \Lambda\ \sim\ \mu\,e^{-1/(2\beta_0g^2)} ,
\label{eq:esssing}
\end{equation}
an essential singularity in the coupling. Every derivative of $e^{-1/g^2}$ vanishes at $g=0$, so
the perturbative gap is identically zero at every order, and analytic continuation in $g^2$ gives
$0$, never the physical value at $g^2=O(1)$. This is a sharper obstruction than the weak-$*$ one
in the following sense: it rules out the \emph{entire analytic-transfer class at once} ---
perturbation theory, Borel resummation, $1/N$, small-torus expansion, and any ``solve exactly
then continue'' scheme. It says: \emph{no procedure analytic in the coupling can manufacture the
nonperturbative scale.} The statement is a theorem about analyticity (the class of functions with
an essential singularity at the origin), not a statement that the gap is uncomputable.

The essential singularity is also the reason the resume/transseries route is blocked at the
level of the condensate: the $u=2$ renormalon ambiguity of the $a^{-4}$ coefficient is the same
size as the coefficient, so the condensate's own ambiguity straddles zero
(\code{ROOF-CONDENSATE-FLOOR.md}). Any floor that passes through the condensate inherits the
ambiguity.

\subsection{The consistency, universality, SUSY and entanglement no-gos}

A wave of ``impossible bases'' tested whether a radically different starting point could bypass
the wall. All eight failed, each for a sharp reason. The table records the outcomes; the
underlying statements are theorems about the relevant structures.

\begin{table}[t]
\centering
\footnotesize
\begin{tabularx}{\textwidth}{@{}>{\raggedright\arraybackslash}p{2.5cm}
  >{\raggedright\arraybackslash}X@{}}
\toprule
\textbf{Base (file)} & \textbf{Why it fails}\\
\midrule
Analytic continuation (\code{IMPOSSIBLE-ANALYTIC}) & \LT{} the gap is an essential singularity
  $e^{-1/(2\beta_0g^2)}$; every analytic-in-$g^2$ method is blind by construction\\
Contradiction (\code{IMPOSSIBLE-CONTRADICTION}) & \LT{} no contradiction obtained; the sharpest
  route needs ``gapless $\Rightarrow$ CFT'' (folklore, not a theorem); it isolates $(\star)$\\
Global consistency (\code{IMPOSSIBLE-CONSISTENCY}) & \LT{} free massless YM satisfies all
  consistency axioms (unitary, Herglotz, crossing, $a_{UV}>a_{IR}$); no 0-form 't Hooft anomaly
  at $\theta=0$; blind\\
Universality (\code{IMPOSSIBLE-UNIVERSALITY}) & \LT{} zero-$T$ crossover is non-universal
  (smooth, no exponents); the one real match (finite-$T$ Ising) is 1-form/center, off-target\\
SUSY + decoupling (\code{IMPOSSIBLE-SUSY-DECOUPLE}) & \LT{} holomorphy computes the condensate
  ($\langle\lambda\lambda\rangle=N\Lambda^3$), not the $0^{++}$ mass; exactness is a symmetry
  pure YM lacks\\
New-space index (\code{IMPOSSIBLE-INDEX}) & \LT{} every index is a homotopy invariant, hence
  maximally continuous; the atom is the maximally discontinuous point\\
Invertibility by degree (\code{IMPOSSIBLE-INVERTIBILITY}) & \LT{} the degree is trivial
  $0/+1$ and unprotected (drops to $0$ at any level crossing $=$ rigidity)\\
Entanglement gap (\code{IMPOSSIBLE-ENTANGLEMENT}) & \LT{} the entanglement gap is the edge/1-form
  spectrum, Schur-orthogonal to $0^{++}$; $\rho_A=\operatorname{Tr}_{A^c}\mu$ is a functional
  of $\mu$\\
\bottomrule
\end{tabularx}
\caption{The eight ``impossible bases'' and why each fails.}
\label{tab:impossible}
\end{table}

The common lesson is the same as in every other wave: the failure is not a lack of cleverness in
the base, but a structural fact about what the base can see. The analytic bases are blind to the
essential singularity; the algebraic bases are rigid; the consistency bases are satisfied by the
free field; the entanglement base is in the wrong channel.

\subsection{The wrong-half, wrong-direction, wrong-regime umbrella}

Three structural failure modes recur across families and are worth naming as no-gos of method:

\begin{itemize}
\item \textbf{Wrong half (UV $\perp$ gap).} L3, H2-global, gradient flow, Dirichlet comparison,
  concentration (Herbst), heat kernel, validated numerics, step scaling, effective string,
  $S^3$ --- these control the UV, a fast rate, or a finite box, but never the IR spectral edge.
\item \textbf{Wrong direction (upper bounds).} Bogoliubov, Hylleraas, finite-size scaling,
  Temple/Weinstein, GEVP --- variational methods give \textbf{upper} bounds on eigenvalues; a
  lower bound needs a threshold oracle, and the oracle is circular.
\item \textbf{Wrong regime (strong coupling only).} Rigorous large-$N$ (Chatterjee), SZZ, random
  surfaces/CPS, center vortices --- these stop at $\beta<\beta_0$ or
  $\lambda_{\mathrm{tH}}<1/24$ and cannot reach $\beta\approx2.5$.
\end{itemize}

\subsection{Summary of no-gos}

\begin{table}[t]
\centering
\footnotesize
\begin{tabularx}{\textwidth}{@{}>{\raggedright\arraybackslash}p{2.8cm}
  >{\raggedright\arraybackslash}p{4.2cm}>{\raggedright\arraybackslash}X@{}}
\toprule
\textbf{No-go} & \textbf{Class it kills} & \textbf{Precise reason}\\
\midrule
Schur-convex & positivity/analyticity-only & free field has $f_1^2=0$; only upper bounds\\
Goldstone & positivity-forced atoms & the only such theorem is at zero mass; BPS wrong channel\\
Periodic & convexification & periodic concave $=$ constant; Savvidy is real\\
Nested cone & second-cone construction & every cone contains the positive cone $\to$ swallows
  free field\\
Rigidity & index/topological & discrete locally-constant index $=$ free value $=0$\\
Feshbach angle & block-elimination & $\Sigma\succeq A$; $\Sigma>0\iff A>0=$ gap\\
Krein & UV spectral invariants & count/sign $\ne$ IR min; split fails at $k=0$\\
Order parameter (B2) & order-parameter/defect & zero-$T$ pure YM is a single phase; nothing to
  condense\\
Markov (BE) & semigroup/flow functionals & homogeneous, zero at equilibrium, or sign-indefinite\\
Analytic transfer & perturbation/Borel/$1/N$ & essential singularity $e^{-1/(2\beta_0g^2)}$\\
Consistency/SUSY/entanglement & global/algebraic bases & satisfied by the free field, rigid, or
  wrong channel\\
Wrong-half/direction/regime & method families & control the wrong scale, sign, or coupling
  region\\
\bottomrule
\end{tabularx}
\caption{Summary of no-gos.}
\label{tab:nogos}
\end{table}

The no-go list is the campaign's negative ledger. It is deliberately long, because a failed
route that is not diagnosed is a route that will be retried. Each row names the class, the reason,
and the file where the reason is proved. The positive ledger is Chapter~\ref{sec:theorems}; the
map that connects them is the subject of the next chapter.
